\section{Appendix}
\label{sec:appendix}

\subsection{DTSE Theory-to-Parameter Mapping}
\label{app:dtse_theory_mapping}
This section anchors the conceptual-to-parameter contract used in Chapter~\ref{sec:methodology}. DTSE stress constructs are mapped to explicit parameter families: Reward Addiction (emission-side controls and burn-to-mint coupling), Subsidy Gap (provider cost and fiat-demand baselines), and Speculative Fragility (price-shock channels and churn sensitivity rules).

\subsection{DTSE Configuration Register}
\label{app:dtse_config_register}
The full parameter register (ranges, distributions, fixed constants, and profile-specific overrides) is maintained as appendix-grade artifacts linked to the DTSE run manifest:
\begin{itemize}
    \item \path{exports/dtse/data/dtse_run_metadata.json}
    \item \path{exports/dtse/tables/dtse_profile_set.csv}
\end{itemize}
Parameter values are treated as modeled assumptions for comparative experiments, not empirical estimates.

\subsection{DTSE Scenario Grid}
\label{app:dtse_scenario_grid}
Table~\ref{tab:app_dtse_scenario_grid} summarizes the frozen DTSE scenario grid used in Chapter~\ref{sec:simulation_results}.
\begin{table}[ht]
    \centering
    \small
    \renewcommand{\arraystretch}{1.2}
    \begin{tabularx}{\textwidth}{@{}>{\raggedright\arraybackslash}p{0.30\textwidth}>{\raggedright\arraybackslash}p{0.20\textwidth}X@{}}
        \toprule
        \textbf{Scenario ID} & \textbf{Label} & \textbf{Description} \\
        \midrule
        \path{baseline_neutral} & Baseline Neutral & Scenario disabled; neutral macro and stable demand for reference trajectories. \\
        \path{demand_contraction} & Demand Contraction & Exogenous usage contraction with higher volatility and bearish background. \\
        \path{liquidity_shock} & Liquidity Shock & Discrete unlock-style sell pressure event with thin liquidity depth. \\
        \path{competitive_yield_pressure} & Competitive Yield Pressure & External yield opportunity raises modeled switching pressure. \\
        \path{provider_cost_inflation} & Provider Cost Inflation & Exogenous increase in provider operating costs under neutral macro. \\
        \bottomrule
    \end{tabularx}
    \caption{DTSE scenario grid summary (frozen Chapter 8 bundle).}
    \label{tab:app_dtse_scenario_grid}
\end{table}
Full grid in \path{exports/dtse/tables/dtse_scenario_grid.csv}.

\subsection{DTSE Scenario-Formula Map (Audit View)}
\label{app:dtse_scenario_formula_map}
Table~\ref{tab:app_dtse_scenario_formula_map} provides a compact audit map linking each Chapter~\ref{sec:simulation_results} scenario to its rationale, core calculation rule, frozen configuration values, and implementation anchors.
\begin{table}[ht]
    \centering
    \footnotesize
    \renewcommand{\arraystretch}{1.2}
    \setlength{\tabcolsep}{3pt}
    \begin{tabularx}{\textwidth}{@{}>{\raggedright\arraybackslash}p{0.17\textwidth}>{\raggedright\arraybackslash}p{0.20\textwidth}>{\raggedright\arraybackslash}p{0.30\textwidth}>{\raggedright\arraybackslash}p{0.17\textwidth}>{\raggedright\arraybackslash}X@{}}
        \toprule
        \textbf{Scenario ID} & \textbf{Why Included} & \textbf{Core Rule (DTSE)} & \textbf{Frozen Inputs / Thresholds} & \textbf{Code / Artifact Anchor} \\
        \midrule
        Baseline neutral &
        Neutral reference path for deviation timing and magnitude comparisons. &
        \(D_t^{\mathrm{BL}}=B_{\mathrm{BL}}(1+\sigma_{\mathrm{BL}}\varepsilon_t)\) with stable regime settings. &
        \path{scenario_id = baseline_neutral}; no scalar threshold crossing required for baseline run definition &
        Export script; scenario grid CSV \\

        Demand contraction &
        Test reward--demand decoupling under exogenous usage deterioration. &
        \(D_t^{\mathrm{DC}}=\max\!\left(0,B_{\mathrm{DC}}f_t(1+\sigma_{\mathrm{DC}}\varepsilon_t)\right)\), where \(f_t\) is the scenario decay function used by the stress generator. &
        \path{scenario_id = demand_contraction}; \path{signal_utilisation_drop_threshold_pct = 0.10} &
        Demand generator; threshold registry CSV \\

        Liquidity shock &
        Test fast transmission from token-liquidity dislocation to provider stress fields. &
        Week-\(20\) unlock sell pressure: \(\Delta Q_{t^*}=S_{t^*}\phi\), then \(P'=k/(X+\Delta Q)^2\). &
        \path{scenario_liquidity_shock_event_week = 20}; \path{scenario_liquidity_shock_investor_sell_pct = 0.35}; \path{signal_price_drop_threshold_pct = 0.10} &
        Simulation engine; threshold registry CSV \\
        \bottomrule
    \end{tabularx}
    \caption{DTSE scenario-formula map for Chapter~\ref{sec:simulation_results} (frozen Chapter 8 bundle).}
    \label{tab:app_dtse_scenario_formula_map}
\end{table}

\begin{table}[ht]
    \ContinuedFloat
    \centering
    \footnotesize
    \renewcommand{\arraystretch}{1.2}
    \setlength{\tabcolsep}{3pt}
    \begin{tabularx}{\textwidth}{@{}>{\raggedright\arraybackslash}p{0.17\textwidth}>{\raggedright\arraybackslash}p{0.20\textwidth}>{\raggedright\arraybackslash}p{0.30\textwidth}>{\raggedright\arraybackslash}p{0.17\textwidth}>{\raggedright\arraybackslash}X@{}}
        \toprule
        \textbf{Scenario ID} & \textbf{Why Included} & \textbf{Core Rule (DTSE)} & \textbf{Frozen Inputs / Thresholds} & \textbf{Code / Artifact Anchor} \\
        \midrule

        Competitive yield pressure &
        Isolate external-opportunity churn without requiring demand collapse. &
        Provider switch pressure \(p_{\mathrm{base}}=0.1\,y^{1.5}\), then type-adjusted probabilities (Provider Type A/B) with bounded noise. &
        \path{scenario_competitive_yield_scalar = 1.5}; \path{signal_providers_drop_threshold_pct = 0.05}; \path{signal_churn_increase_threshold_pct = 0.20} &
        Simulation engine; threshold registry CSV \\

        Provider cost inflation &
        Test Subsidy-Gap exposure through exogenous operating-cost increase. &
        \(c_{\mathrm{base}}^{\mathrm{PCI}}=1.35\,c_{\mathrm{base}}^{\mathrm{BL}}\), propagated through heterogeneous provider-cost rule. &
        \path{scenario_provider_cost_multiplier = 1.35}; \path{signal_profit_drop_threshold_pct = 0.10}; \path{signal_providers_drop_threshold_pct = 0.05} &
        Export script; simulation engine \\
        \bottomrule
    \end{tabularx}
    \caption[]{DTSE scenario-formula map for Chapter~\ref{sec:simulation_results} (continued).}
\end{table}
In this audit map, values listed in the ``Frozen Inputs / Thresholds'' column are scalar fields exported in the frozen scenario grid or threshold registry. Symbolic constants in the ``Core Rule'' column are equation-level notation and are not all exported as standalone scalar registry fields.

\subsection{DTSE Threshold Definitions}
\label{app:dtse_threshold_table}
Table~\ref{tab:app_dtse_threshold_registry_excerpt} provides a concise excerpt of the exported scalar threshold registry used to interpret Chapter~\ref{sec:simulation_results}.
\begin{table}[ht]
    \centering
    \footnotesize
    \renewcommand{\arraystretch}{1.2}
    \setlength{\tabcolsep}{3pt}
    \begin{tabularx}{\textwidth}{@{}>{\raggedright\arraybackslash}p{0.34\textwidth}>{\raggedright\arraybackslash}p{0.08\textwidth}>{\raggedright\arraybackslash}p{0.13\textwidth}>{\raggedright\arraybackslash}p{0.15\textwidth}>{\raggedright\arraybackslash}X@{}}
        \toprule
        \textbf{Threshold ID} & \textbf{Value} & \textbf{Unit} & \textbf{Scope} & \textbf{Use in Chapter 8} \\
        \midrule
        \path{fm_profitability_negative_profit_threshold} & \texttt{0} & \path{usd_per_week} & \path{chapter8_failure_mode} & Profitability-induced churn criterion. \\
        \path{fm_profitability_sustained_window_weeks} & \texttt{3} & \path{weeks} & \path{chapter8_failure_mode} & Sustained window for \(N\) in failure-mode prose. \\
        \path{scenario_liquidity_shock_event_week} & \texttt{20} & \path{week_index} & \path{scenario_config} & Event marker used in liquidity-shock reporting. \\
        \path{scenario_liquidity_shock_investor_sell_pct} & \texttt{0.35} & \path{fraction} & \path{scenario_config} & Unlock-style sell-pressure scalar in liquidity-shock scenario. \\
        \path{scenario_competitive_yield_scalar} & \texttt{1.50} & \path{relative_yield_multiple} & \path{scenario_config} & External-yield pressure scalar used in competitive-yield scenario. \\
        \path{scenario_provider_cost_multiplier} & \texttt{1.35} & \path{multiplier} & \path{scenario_config} & Exogenous provider-cost multiplier in cost-inflation scenario. \\
        \path{signal_price_drop_threshold_pct} & \texttt{0.10} & \path{fraction} & \path{signal_detection} & Price-signal detection threshold versus baseline. \\
        \path{signal_utilisation_drop_threshold_pct} & \texttt{0.10} & \path{fraction} & \path{signal_detection} & Utilization-signal detection threshold versus baseline. \\
        \path{signal_profit_drop_threshold_pct} & \texttt{0.10} & \path{fraction} & \path{signal_detection} & Profit-signal detection threshold versus baseline. \\
        \path{signal_providers_drop_threshold_pct} & \texttt{0.05} & \path{fraction} & \path{signal_detection} & Provider-count signal threshold versus baseline. \\
        \path{signal_churn_increase_threshold_pct} & \texttt{0.20} & \path{fraction} & \path{signal_detection} & Churn-increase threshold versus baseline. \\
        \path{profile_ono_v3_calibrated_churn_threshold} & \texttt{10} & \path{usd/week_profit} & \path{profile_param} & ONO provider exit sensitivity scalar. \\
        \bottomrule
    \end{tabularx}
    \caption{Excerpt of scalar threshold registry for frozen Chapter 8 DTSE artifacts.}
    \label{tab:app_dtse_threshold_registry_excerpt}
\end{table}
Full registry in \path{exports/dtse/tables/dtse_threshold_registry.csv}.

\subsection{Override Ledger (Profile-Specific Parameters)}
\label{app:dtse_override_ledger}
This table documents profile-specific overrides and global heterogeneity rules used in the DTSE runs reported in Chapter~\ref{sec:simulation_results}. All entries listed here are modeled assumptions (including calibration choices) rather than empirical estimates of real-world protocol behavior. Each entry is traceable either to frozen exported artifacts (scenario grid, threshold registry, run manifest) or to explicit simulator-rule definitions when the rule is not exported as a standalone scalar field. The table is provided for transparency and reproducibility and does not imply causal attribution or performance claims.
Table~\ref{tab:app_dtse_override_ledger} is the audit ledger for these profile-specific and global rule overrides.

\begin{table}[ht]
    \centering
    \footnotesize
    \renewcommand{\arraystretch}{1.2}
    \setlength{\tabcolsep}{3pt}
    \begin{tabularx}{\textwidth}{@{}>{\raggedright\arraybackslash}p{0.29\textwidth}>{\raggedright\arraybackslash}p{0.09\textwidth}>{\raggedright\arraybackslash}p{0.12\textwidth}>{\raggedright\arraybackslash}X>{\raggedright\arraybackslash}p{0.16\textwidth}@{}}
        \toprule
        \textbf{Key / ID} & \textbf{Scope} & \textbf{Class} & \textbf{Purpose / Evidence Status} & \textbf{Artifact Anchor} \\
        \midrule
        \path{profile_ono_v3_calibrated_churn_threshold} &
        ONO &
        \path{profile_param} &
        Profile-specific churn sensitivity threshold used to parameterize provider exit response under stress; modeled calibration assumption (not an empirical estimate). &
        Threshold registry; run manifest \\

        \path{profile_ono_v3_calibrated_max_provider_churn_rate} &
        ONO &
        \path{model_param} &
        Profile-specific upper bound on provider churn dynamics used for stability and regime control in simulation; modeled assumption for the ONO profile. &
        Threshold registry; run manifest \\

        \path{profile_helium_bme_v1_churn_threshold} &
        Helium &
        \path{profile_param} &
        Profile-specific churn sensitivity threshold used to parameterize provider exit response under stress for the Helium comparator profile; modeled assumption within DTSE. &
        Threshold registry; run manifest \\
        \bottomrule
    \end{tabularx}
    \caption{Override ledger for DTSE Chapter 8 runs (profile-specific parameters and global heterogeneity rules).}
    \label{tab:app_dtse_override_ledger}
\end{table}

\begin{table}[ht]
    \ContinuedFloat
    \centering
    \footnotesize
    \renewcommand{\arraystretch}{1.2}
    \setlength{\tabcolsep}{3pt}
    \begin{tabularx}{\textwidth}{@{}>{\raggedright\arraybackslash}p{0.29\textwidth}>{\raggedright\arraybackslash}p{0.09\textwidth}>{\raggedright\arraybackslash}p{0.12\textwidth}>{\raggedright\arraybackslash}X>{\raggedright\arraybackslash}p{0.16\textwidth}@{}}
        \toprule
        \textbf{Key / ID} & \textbf{Scope} & \textbf{Class} & \textbf{Purpose / Evidence Status} & \textbf{Artifact Anchor} \\
        \midrule

        \path{profile_geodnet_v1_churn_threshold} &
        Geodnet &
        \path{profile_param} &
        Profile-specific churn sensitivity threshold used to parameterize provider exit response under stress for the Geodnet comparator profile; modeled assumption within DTSE. &
        Threshold registry; run manifest \\

        \path{competitive_yield_provider_type_adjustment_rule} &
        Global &
        \path{scenario_logic} &
        Global heterogeneity rule within the competitive-yield stress mechanism that applies differential switching sensitivity across provider types; scenario-logic construct anchored to the simulator specification rather than protocol facts. This rule is code-level and not exported as a standalone scalar registry field. &
        Scenario formula map; simulation engine \\
        \bottomrule
    \end{tabularx}
    \caption[]{Override ledger for DTSE Chapter 8 runs (continued).}
\end{table}

\subsection{DTSE Run Manifest}
\label{app:dtse_run_manifest}
Table~\ref{tab:app_dtse_manifest_excerpt} provides a concise excerpt of the frozen run manifest used for Chapter 8 reporting.
\begin{table}[ht]
    \centering
    \small
    \renewcommand{\arraystretch}{1.2}
    \begin{tabularx}{\textwidth}{@{}>{\raggedright\arraybackslash}p{0.30\textwidth}X@{}}
        \toprule
        \textbf{Manifest Field} & \textbf{Value} \\
        \midrule
        DTSE repo identifier & \path{DePin-Stress-Test} (external repository) \\
        DTSE commit hash & \path{0f7b69d0b5776137bddcea6ed9fc642259764792} \\
        Run timestamp (UTC) & \texttt{2026-02-24T10:21:41.956Z} \\
        Scenario grid ID & \texttt{dtse\_ch8\_grid\_v1} \\
        Seed policy & \texttt{seed = MASTER\_SEED + sim\_index}; \texttt{master\_seed = 20260224} \\
        Horizon / simulations & \texttt{horizon\_weeks = 52}; \texttt{n\_sims = 60} \\
        Mechanism profile IDs & \shortstack[l]{\path{ono_v3_calibrated},\\\path{helium_bme_v1},\\\path{hivemapper_v1},\\\path{grass_v1},\\\path{geodnet_v1}} \\
        Scenario IDs & \shortstack[l]{\path{baseline_neutral},\\\path{demand_contraction},\\\path{liquidity_shock},\\\path{competitive_yield_pressure},\\\path{provider_cost_inflation}} \\
        \bottomrule
    \end{tabularx}
    \caption{Run manifest excerpt for frozen Chapter 8 DTSE artifacts.}
    \label{tab:app_dtse_manifest_excerpt}
\end{table}
Full JSON manifest (including artifact SHA256 hashes) at \path{exports/dtse/run_manifest_2026-02-24.json}.

\subsection{DTSE Metric Computation Notes}
\label{app:dtse_metric_computation}
Metric computation rules provide the exact formulas and aggregation conventions used in reporting: per-run timeseries definitions, cross-run summaries (median/IQR), and derived diagnostic signatures used in failure-mode classification.
\begin{itemize}
    \item \path{exports/dtse/tables/dtse_weekly_summary.csv}
    \item \path{exports/dtse/tables/dtse_final_summary.csv}
\end{itemize}

\subsection{Empirical Metric Applicability Catalogue}
\label{app:empirical_metric_applicability}
This section documents the metric-applicability contract used in Chapter~\ref{sec:empirical_analysis}. It governs only empirical event-window interpretation and does not modify the pre-specified DTSE reporting set in Chapter~\ref{sec:simulation_results}.
Table~\ref{tab:app_empirical_metric_catalogue} reports the cross-project empirical applicability catalogue, and Table~\ref{tab:app_onocoy_metric_applicability} reports the Onocoy-specific applicability profile.

\begin{table}[ht]
    \centering
    \footnotesize
    \renewcommand{\arraystretch}{1.2}
    \setlength{\tabcolsep}{4pt}
    \begin{tabularx}{\textwidth}{@{}>{\raggedright\arraybackslash}p{0.20\textwidth}>{\raggedright\arraybackslash}p{0.16\textwidth}>{\raggedright\arraybackslash}p{0.16\textwidth}>{\raggedright\arraybackslash}p{0.13\textwidth}>{\raggedright\arraybackslash}X@{}}
        \toprule
        \textbf{Metric} & \textbf{Mechanism Compatibility} & \textbf{Empirical Data Requirement} & \textbf{Early-Stage Status} & \textbf{Interpretive Rule} \\
        \midrule
        Burn-to-Mint Ratio & Native in BME; adapted in capped-supply reward-distribution designs & Observable burn/sink and reward-distribution fields & Partial & Directional interpretation only when adaptation is used; avoid direct magnitude equivalence across mechanism classes. \\
        30-Day Retention/Churn & Mechanism-agnostic & Time-indexed active-provider observations across stress windows & Partial & Directional under short histories; stronger inference when stress-window depth is sufficient. \\
        Revenue per Node & Mechanism-agnostic & Stable revenue denominator plus comparable provider counts & N/R & Mark N/R when revenue depth is insufficient or economically unstable (NR3). \\
        Token Turnover & Mechanism-agnostic & Sustained market liquidity and turnover observability & N/R & Mark N/R when liquidity is thin and turnover is microstructure-dominated (NR3). \\
        FDV / Annualized Revenue & Mechanism-agnostic & Reliable market-cap proxy and stable annualized revenue denominator & N/R & Exclude from ranking when denominator instability drives extreme ratios (NR3). \\
        \bottomrule
    \end{tabularx}
    \caption{Empirical metric-applicability catalogue used in Chapter~\ref{sec:empirical_analysis}.}
    \label{tab:app_empirical_metric_catalogue}
\end{table}

\begin{table}[ht]
    \centering
    \footnotesize
    \renewcommand{\arraystretch}{1.2}
    \setlength{\tabcolsep}{4pt}
    \begin{tabularx}{\textwidth}{@{}>{\raggedright\arraybackslash}p{0.24\textwidth}>{\raggedright\arraybackslash}p{0.12\textwidth}>{\raggedright\arraybackslash}p{0.18\textwidth}X@{}}
        \toprule
        \textbf{Onocoy Metric} & \textbf{Status} & \textbf{Constraint Type} & \textbf{Interpretive Handling in Chapter~\ref{sec:empirical_analysis}} \\
        \midrule
        Burn-to-Mint (adapted) & Partial & Mechanism adaptation & Used directionally as burn-to-distribution trend under capped-supply architecture. \\
        Retention / Churn & Partial & Temporal depth & Used as primary participation signal with early-stage caveat. \\
        Revenue per Node & N/R & NR3 & Excluded from cross-project ranking in current empirical windows. \\
        Token Turnover & N/R & NR3 & Excluded until sustained ONO liquidity supports stable interpretation. \\
        FDV / Annualized Revenue & N/R & NR3 & Excluded where denominator instability dominates metric value. \\
        \bottomrule
    \end{tabularx}
    \caption{Onocoy empirical metric-applicability profile for current stress windows.}
    \label{tab:app_onocoy_metric_applicability}
\end{table}
